\section{задание: Решение логических задач}
\subsection{исходные данные}
\indent Тип задачи: \textbf{Логика (Переливания)} \\
\indent Описание задачи (кратко): \textbf{Задача о сосудах: Отмерить ровно 4 литра воды, имея только два сосуда емкостью 3 и 5 литров и неограниченный доступ к крану.} \\
\indent Сложность: \textbf{Низкая} \\
\newline
\noindent Изменённая задача: \\
\indent \textbf{Задача о ложках: Отмерить ровно 6 грамм сахара, имея только две ложки объёмом 7 и 10 грамм и неограниченный сахар.} \\
\subsection{базовый промпт}
\begin{verbatim}
					Реши  эту  задачу: {ЗАДАЧА}
\end{verbatim}
<<Наивный>> запрос из примера.
\subsection{журнал промпт-инжиниринга}
\noindent Промпт v1:
\begin{verbatim}
					Реши  эту  задачу: Задача о ложках: Отмерить ровно 6 грамм
					сахара, имея только две ложки объёмом 7 и 10 грамм и
					неограниченный сахар.
\end{verbatim}
\newpage
\noindent Результат (YandexGPT):
\begin{figure}[H]
	\centering
	\includegraphics[width=15cm]{z7/a1.png}
\end{figure}
\noindent Анализ ошибок: модель без дополнительных запросов выдала ответ с рассуждением.\\
Корректировка: не требуется. \\
\newpage
\noindent Результат (DeepSeek):
\begin{figure}[H]
	\centering
	\includegraphics[width=13cm]{z7/b1.png}
\end{figure}
\newpage
\noindent Анализ ошибок: ответ слишком длинный. \\
Корректировка: попросить модель писать короче. \\
\newline
\noindent Промпт v2:
\begin{verbatim}
					Давай  рассуждать  шаг  за  шагом, кратко, но ёмко.
					Реши  эту  задачу: Задача о ложках: Отмерить ровно 6 грамм
					сахара, имея только две ложки объёмом 7 и 10 грамм и
					неограниченный сахар.
\end{verbatim}
\newpage
\noindent Результат (YandexGPT):
\begin{figure}[H]
	\centering
	\includegraphics[width=15cm]{z7/c1.png}
\end{figure}
\noindent Анализ ошибок: модель выдала такой же ответ. \\
Корректировка: не требуется. \\
\newpage
\noindent Результат (DeepSeek):
\begin{figure}[H]
	\centering
	
\includegraphics[width=15cm]{z7/d1.png}
\end{figure}
\noindent Анализ ошибок: модель выдала расписанный ответ. \\
Корректировка: не требуется. \\
\newpage
\subsection{Сравнение моделей:}
\begin{table}[h!]
	\centering
	\begin{tabular}{|l|p{5cm}|p{5cm}|}
		\hline
		\textbf{Критерий} & \textbf{YandexGPT} & \textbf{DeepSeek} \\
		\hline
		Логическая связность & 5 & 5 \\
		\hline
		Корректность финального ответа & 5 & 5 \\
		\hline
		\textbf{Итог} & 5 & 5 \\
		\hline
	\end{tabular}
\end{table}
Обе модели одинакого хорошо решают эту простую задачу на своей базовой версии. DeepSeek расписывает подробнее так как это его особенность.
\subsection{Финальный промпт:}
\begin{verbatim}
					Давай  рассуждать  шаг  за  шагом, кратко, но ёмко.
					Реши  эту  задачу: {ЗАДАЧА}
\end{verbatim}
\subsection{Вывод по заданию:}
\indent Это задание демонстрирует что самые доступные модели хорошо справлятся с простыми логическими задачами, без дополнительных указаний. Так как многие из них имееют встроенный движок для <<обдумывания>>. 